\section{马尔可夫决策过程}
\label{sec:mdp}

\glsfirst{mdp}（马尔可夫决策过程）\cite{puterman1994markov} 模型可视为一种 \gls{mc}，其转移受动作执行的影响。通常，智能体负责选择动作，从而能与 \gls{mdp} 进行交互。此外，转移过程关联着一个奖励函数，为智能体提供反馈。

形式化地，\gls{mdp} 定义为一个五元组 $\markovdecisionprocess = (\statespace,\actionspace, \transitionfunction, \rewardfunction, \initialstatedistribution)$：$\statespace$ 和 $\actionspace$ 分别是状态和动作的可测集；$\transitionfunction(s'\vert s,a)$ 是通过执行动作 $a$ 从状态 $s$ 转移到状态 $s'$ 的概率；$R(s,a)$ 是一个随机变量，表示智能体在施加动作 $a$ 的状态 $s$ 转移过程中获得的奖励；$\initialstatedistribution$ 是智能体初始状态的分布。记 $r(s,a)=\expectedvalue[R(s,a)]$，并在本文中做出以下假设。

\begin{ass}[有界奖励]
    奖励是有界的，若 $\forall (s,a)\in\statespace\times\actionspace$，
    \begin{align*}
        0\le r(s,a)\le \maximumreward.
    \end{align*}
\end{ass}

%We will also use the notation $S_t$ or $A_t$ to refer to the random variable holding the value of the state or the action after $t$ steps in the \gls{mdp}.

\subsection{策略}
\label{subsec:policies}
智能体根据其与环境的交互历史 $h_t=(s_0,a_0,\dots,s_{t-1},a_{t-1},s_t)$ 选择动作的方式称为\emph{策略}。若记 $\mathcal{H}_t$ 为 $t$ 时刻所有可能历史的集合，则可形式化地将智能体策略定义为从历史集到动作空间上概率集的映射：
\begin{align*}
    \pi:\historyspace_t\rightarrow\setofprobability\left(\actionspace\right).
\end{align*}
本文沿用 \cite{puterman1994markov} 的符号，记所有策略的集合为 $\Pi^{\text{HR}}$，其中 "$R$" 表示随机化，"$H$" 表示历史依赖。

$\Pi^{\text{HR}}$ 的不同子集因其性质而值得关注。马尔可夫策略子集考虑仅依赖于最近观测状态的策略。这是一个便利的性质，因为它避免了策略输入随时间增长维度。已有研究表明，对于后文将介绍的目标函数，该子集中包含最优策略 \cite[Theorem~6.2.10 and Theorem~8.1.2]{puterman1994markov}。
%\pl{valid only for discrete state space? see https://cs.stackexchange.com/questions/93617/existence-of-optimal-policy-for-infinite-state-mdps}
使用上标 "$M$" 表示马尔可夫策略。当策略随时间固定（尽管可能随机）时，称其为平稳策略，并使用上标 "$S$"。最后，当 $\actionspace$ 上的概率测度退化、将所有质量集中于单一动作时，策略是确定性的。该子集用字母 "$D$" 标记。综上所述，这些集合的包含关系如下 \cite[Section 2.1.5]{puterman1994markov}：
\begin{align*}
    \begin{array}{cccc}
         \Pi^{\text{SD}}&\subset\Pi^{\text{SR}}&\subset\Pi^{\text{MR}}&\subset\Pi^{\text{HR}}  \\
         \Pi^{\text{SD}}&\subset\Pi^{\text{MD}}&\subset\Pi^{\text{MR}}&\subset\Pi^{\text{HR}}\\
         \Pi^{\text{SD}}&\subset\Pi^{\text{MD}}&\subset\Pi^{\text{HD}}&\subset\Pi^{\text{HR}}.
    \end{array}
\end{align*}


\subsection{目标}
\label{subsec:objective}
\gls{mdp} 可定义三种主要目标：期望总回报、期望折扣回报和平均奖励 \cite{puterman1994markov}。本文将关注后两类回报。对于某个 \gls{mdp} $\markovdecisionprocess$，令 $\probability\left(S_t=s; \pi,\initialstatedistribution\right)$ 为在策略 $\pi$ 和初始状态分布 $\initialstatedistribution$ 下，表示 $t$ 时刻状态的随机变量 $S_t$ 取值为 $s$ 的概率。对于某个策略 $\policy$，折扣因子为 $\gamma$ 的期望折扣回报定义为
\begin{align}
\label{eq:discounted_reward}
    \expectedreturn[\gamma] = \expectedvalue_{\substack{s_{t+1}\sim p(\cdot\vert s_t,a_t)\\a_t\sim\pi(\cdot\vert s_t)\\s_0\sim \initialstatedistribution}}\left[\sum_{t=1}^H \gamma^t r(s_t,a_t) \right];
\end{align}
其中 $H$ 是可能为无穷大的水平线。而平均奖励性能为\footnote{此处沿用 $H$ 而非更常用的 $T$，以匹配折扣情形的符号约定，并将 $T$ 保留用于表示跨回合的总步数。}
\begin{align*}
    \averagereturnfunction = \lim_{H\rightarrow \infty}\frac{1}{H}\expectedvalue_{\substack{s_{t+1}\sim p(\cdot\vert s_t,a_t)\\a_t\sim\pi(\cdot\vert s_t)\\s_0\sim \initialstatedistribution}}\left[\sum_{t=1}^H r(s_t,a_t) \right].
\end{align*}

这些期望值针对初始状态分布 $d_0$ 和策略诱导的转移分布计算。

\subsection{状态-动作占用度量和分布}
\label{subsec:def_state_distrib}

从理论角度来看，对智能体访问的状态具有某种度量或分布是有用的。实际上，策略的性能可直接从这些量推导得出。首先定义折扣回报准则下的状态-动作占用分布：
\begin{align}
\label{def:discounted_state_distrib}
    \discountedstateoccupancydistribution[\gamma] (s,a)= (1-\gamma)\expectedvalue_{\substack{s_{t+1}\sim p(\cdot\vert s_t,a_t)\\a_t\sim\pi(\cdot\vert s_t)\\s_0\sim \initialstatedistribution}}\left[\sum_{t=1}^H \gamma^t \Ind\left(S_t=s, A_t=a\right)\right].
\end{align}
而平均奖励下的状态-动作占用分布为
\begin{align*}
    \averagestateoccupancydistribution (s,a)= \lim_{H\rightarrow \infty}\frac{1}{H}\expectedvalue_{\substack{s_{t+1}\sim p(\cdot\vert s_t,a_t)\\a_t\sim\pi(\cdot\vert s_t)\\s_0\sim \initialstatedistribution}}\left[\sum_{t=1}^H \Ind\left(S_t=s, A_t=a\right)\right].
\end{align*}
% for the total reward criterion,
% \begin{align*}
%     \totalstateoccupancymeasure[\gamma] = \expectedvalue\left[\sum_{t=1}^\infty  \probability^\policy(s_t=s) \pi(a\vert s)\right].ò
% \end{align*}
从状态-动作分布可推导出状态分布如下：
\begin{align*}
    \discountedstateoccupancydistribution[\gamma] (s)&=\int_\actionspace \discountedstateoccupancydistribution[\gamma] (s,a) \de a,
    \\
    \averagestateoccupancydistribution (s) &= \int_\actionspace\averagestateoccupancydistribution (s,a)\de a.
\end{align*}
注意此处使用了相同的符号，但函数的输入可避免混淆。为简化符号，当目标不会引起混淆时，有时会省略下标 ``$\gamma$'' 或 ``AVG''。

最后，定义折扣情形下的状态-动作占用度量 \cite{laroche2022non}：
\begin{align*}
    \discountedstateoccupancymeasure[\gamma] (s,a) &= \expectedvalue_{\substack{s_{t+1}\sim p(\cdot\vert s_t,a_t)\\a_t\sim\pi(\cdot\vert s_t)\\s_0\sim \initialstatedistribution}}\left[\sum_{t=1}^H \gamma^t \Ind\left(S_t=s, A_t=a\right)\right]
    \\
    &= \frac{\discountedstateoccupancydistribution[\gamma] (s,a)}{1-\gamma},
\end{align*}
以及平均奖励下的对应量：
\begin{align*}
    \averagestateoccupancymeasure[\pi] (s,a) = \lim_{H\rightarrow \infty}\expectedvalue_{\substack{s_{t+1}\sim p(\cdot\vert s_t,a_t)\\a_t\sim\pi(\cdot\vert s_t)\\s_0\sim \initialstatedistribution}}\left[\sum_{t=1}^H \Ind\left(S_t=s, A_t=a\right)\right].
\end{align*}
策略的折扣回报可从状态-动作占用度量推导如下 \cite[Lemma~1]{laroche2022non}：
\begin{align*}
    \expectedreturn[\gamma] &= \int_{\statespace,\actionspace} r(s,a)d\discountedstateoccupancymeasure[\gamma]  (\de s,\de a).
\end{align*}

\subsection{价值函数}
在 \gls{rl} 中，量化策略性能的重要概念是\emph{价值函数}。在折扣情形下，状态-动作价值函数是智能体从某个状态 $s$ 出发、先执行动作 $a$ 而后遵循策略 $\pi$ 所获得的期望折扣奖励总和。其形式为：
\begin{align}
    Q^\pi(s,a) = \expectedvalue_{\substack{s_{t+1}\sim p(\cdot\vert s_t,a_t)\\a_t\sim\pi(\cdot\vert s_t)}}
    \left[\left.\sum_{t=0}^H \gamma^t r(s_t,a_t)\right\vert s_0=s, a_0=a\right].\label{eq:state_action_value_function}
\end{align}
该函数称为\emph{Q函数}（Q-function）。自然也可定义状态价值函数（或简称价值函数），它是智能体从某个状态 $s$ 出发并立即遵循策略 $\pi$ 所获得的期望折扣奖励总和。状态价值函数由 Q 函数定义为：
\begin{align*}
    V^\pi(s)=\expectedvalue_{a\sim\pi(\cdot\vert s)}[Q^\pi(s,a)].
\end{align*}

用于比较动作效果的一个有用概念是优势函数 $\advantagefunction$。它量化了仅在下一步偏离给定策略所获得的优势。形式化地，
\begin{align}
\label{eq:advantage_function}
    \advantagefunction(s,a) = Q^\pi(s,a) - V^\pi(s)
\end{align}

在 \gls{rl} 中，关注的是状态-动作价值函数和状态价值函数所能达到的最大值。最优状态价值函数和最优状态-动作价值函数为：
\begin{align*}
    \optimalstatevaluefunction (s)&= \sup_{\pi\in\Pi^{\text{HR}}}V^\pi(s),
    \\
    \optimalstateactionvaluefunction (s,a)&= \sup_{\pi\in\Pi^{\text{HR}}}Q^\pi(s,a).
\end{align*}
达到这些值的策略 $\pi\in\Pi^{\text{HR}}$ 记为 $\optimalpolicy$。注意 $\Pi^{\text{SR}}$ 中存在最优策略 \cite[Theorem~6.2.10]{puterman1994markov}，因此可将最优策略的搜索限制在该集合中。

\subsection{Bellman方程}
在 \glspl{mdp} 中寻找最优价值函数的核心概念是\emph{动态规划} \cite{bellman1954theory}，其思想是将问题分解为更小的子问题。应用于 \gls{rl} 时，计算任意状态价值函数的问题被递归地分解为每一步的价值函数关于未来价值函数的计算。这得益于\emph{Bellman期望算子} \cite{bellman1966dynamic}。

\begin{defi}[Bellman期望算子]
    考虑一个 \gls{mdp} 和策略 $\pi\in\Pi^{\text{SR}}$，令 $\mathcal{B}(\statespace)$ 为 $\statespace$ 上有界可测函数的集合。状态价值函数的 Bellman 期望算子 $\bellmanoperator_V:\mathcal{B}(\statespace)\mapsto\mathcal{B}(\statespace)$ 对 $f\in\mathcal{B}$ 和 $s\in\statespace$ 定义如下：
    \begin{align}
    \label{eq:bellman_operator_V}
        (\bellmanoperator_V f)(s) = \int_{\actionspace}\pi(\de a\vert s)\left(r(s,a) + \gamma  \int_\statespace  \transitionfunction(\de s'\vert s,a) f(s')\right).
    \end{align}
    类似地，状态-动作价值函数的 Bellman 期望算子 $\bellmanoperator_Q:\mathcal{B}(\statespace\times\actionspace)\mapsto\mathcal{B}(\statespace\times\actionspace)$ 对 $f\in\mathcal{B}$、$s\in\statespace$ 和 $a\in\actionspace$ 定义如下：
    \begin{align}
    \label{eq:bellman_operator_Q}
        (\bellmanoperator_Q f)(s,a) = r(s,a) + \gamma  \int_\statespace  \transitionfunction(\de s'\vert s,a)\int_\actionspace \pi(\de a'\vert s')f(s',a').
    \end{align}
\end{defi}

若 $\gamma<1$，这些算子是 $L_\infty$-范数下的 $\gamma$-压缩映射 \cite[Proposition~6.2.5]{puterman1994markov}。回忆 $L_\infty$-范数下的 $\gamma$-压缩映射是指算子 $T:X\mapsto X$ 对 $f,g\in X$ 满足：
\begin{align*}
    \lVert Tf - Tg\rVert_\infty\le\gamma\lVert f-g\rVert_\infty.
\end{align*}
Bellman期望算子的这些性质使其满足Banach不动点定理的条件，意味着递归应用这些算子将收敛到唯一的不动点。Q 函数是 \Cref{eq:bellman_operator_Q} 的不动点，状态价值函数是 \Cref{eq:bellman_operator_V} 的不动点。由此定义\emph{Bellman期望方程}：
\begin{align}
    Q^\pi &= r(s,a) + \gamma  \int_\statespace  \transitionfunction(\de s'\vert s,a)V^\pi(s'),
    \label{eq:bellman_equation_Q}\\
    V^\pi &= \int_\actionspace \pi(\de a\vert s) Q^\pi(s,a).
    \label{eq:bellman_equation_V}
\end{align}

有趣的是，最优状态-动作价值函数和状态价值函数也是类似算子的不动点，这些算子称为\emph{Bellman最优性算子}，定义如下。

\begin{defi}[Bellman最优性算子]
    考虑一个 \gls{mdp}，令 $\mathcal{B}(\statespace)$ 为 $\statespace$ 上有界可测函数的集合。状态价值函数的 Bellman 最优性算子 $\optimalbellmanoperator_V:\mathcal{B}(\statespace)\mapsto\mathcal{B}(\statespace)$ 对 $f\in\mathcal{B}$ 和 $s\in\statespace$ 定义如下：
    \begin{align}
    \label{eq:bellman_optimal_operator_V}
        (\optimalbellmanoperator_V f)(s) = \sup_{a\in\actionspace}\left\{r(s,a) + \gamma  \int_\statespace  \transitionfunction(\de s'\vert s,a) f(s')\right\}.
    \end{align}
    类似地，状态-动作价值函数的 Bellman 期望算子 $\optimalbellmanoperator_Q:\mathcal{B}(\statespace\times\actionspace)\mapsto\mathcal{B}(\statespace\times\actionspace)$ 对 $f\in\mathcal{B}$、$s\in\statespace$ 和 $a\in\actionspace$ 定义如下：
    \begin{align}
    \label{eq:bellman_optimal_operator_Q}
        (\optimalbellmanoperator_Q f)(s,a) = r(s,a) + \gamma  \int_\statespace  \transitionfunction(\de s'\vert s,a)\sup_{a'\in\actionspace} \left\{f(s',a')\right\}.
    \end{align}
\end{defi}

这些算子也是 $\gamma$-压缩映射 \cite[Proposition~6.2.4]{puterman1994markov}，根据Banach不动点定理，分别以最优状态价值函数和最优状态-动作价值函数为不动点 \cite[Proposition~6.2.5]{puterman1994markov}。

类似地，最优价值函数满足以下\emph{Bellman最优性方程}：
\begin{align}
    Q^\star(s,a) &= r(s,a) + \gamma  \int_\statespace  \transitionfunction(\de s'\vert s,a)V^\star(s'),
    \label{eq:bellman_optimal_equation_Q}\\
    V^\star(s) &= \sup_{a\in\actionspace} Q^\star(s,a).
    \label{eq:bellman_optimal_equation_V}
\end{align}


\subsection{光滑马尔可夫决策过程}
\label{subsec:smooth_mdp}

考虑所有 \glspl{mdp} 意味着需要考虑一些非常复杂、难以高效求解的问题。对 \glspl{mdp} 施加假设以缩小范围，可以大幅降低复杂度，同时保持模型的现实性。在 \gls{rl} 中普遍存在的一个假设是关于 \gls{mdp} 的\emph{光滑性}。为定义光滑性，首先回顾Lipschitz性质的概念。

\begin{defi}[Lipschitz连续性]
    考虑 $(X,\distance_X)$ 和 $(Y,\distance_Y)$ 两个度量空间，以及函数 $f:X \rightarrow Y$。若对于 $L>0$ 有
    \begin{align*}
        \forall x,x'\in X,\; \distance_Y(f(x),f(x'))\leq L \distance_X(x,x'),
    \end{align*}
    则称 $f$ 是 $L$-Lipschitz连续的（$L$-LC）。
\end{defi}

在本文中，对于 $X\subset \mathbb{R}^n$ 的集合 $X$，考虑Euclidean距离，定义为 $\distance_{X}(x,x')=\lVert x-x'\rVert_2$。对于概率分布，将考虑 $L_1$-Wasserstein距离\footnote{为简洁起见，本文中将 $L_1$-Wasserstein 距离简称为 Wasserstein 距离。}，其定义如下。

\begin{defi}[$L_1$-Wasserstein距离，\cite{villani2009optimal}]
    设 $\mu,\nu$ 为样本空间 $\Omega$ 上的两个概率测度：
    \begin{align*}
        \wassersteindistance[1](\mu\Vert \nu)=\sup_{\left\Vert f\right\Vert_L\leq 1} \left\vert \int_{\Omega} f(\omega)(\mu-\nu)\; (\de \omega) \right\vert,
    \end{align*}
    其中 $\left\Vert f\right\Vert_L$ 是 $f$ 的Lipschitz半范数：
    \begin{align*}
        \left\Vert f\right\Vert_L=\sup_{x,x'\in X, x\neq x'} \frac{\distance_Y(f(x),f(x'))}{\distance_X(x,x')}.
    \end{align*}
\end{defi}

Lipschitz性质可理解为某种光滑性概念，因为它限制了函数的增长速度。现准备描述 \gls{mdp} 的光滑性。

\begin{defi}[Lipschitz \gls{mdp}，\cite{rachelson2010locality}]
\label{def:lip_mdp}
    若 $\forall(s,a),(s',a')\in\statespace\times\actionspace$ 有
    \begin{align*}
        & \wassersteindistance[1](\transitionfunction(\cdot\vert s,a)\Vert \transitionfunction(\cdot\vert s',a'))\leq L_P \left( \distance_{\statespace}(s,s')  + \distance_{\actionspace}(a,a')\right),
        \\
        & \left\vert r(s,a)-r(s',a') \right\vert \leq L_r \left( \distance_{\statespace}(s,s') + \distance_{\actionspace}(a,a')\right),
    \end{align*}
    则称该 \gls{mdp} 是 $(L_P,L_r)$-LC 的。
\end{defi}

类似的光滑性假设也可施加于策略，如下定义所示。

\begin{defi}[Lipschitz策略]
\label{def:lip_policy}
    一个平稳马尔可夫策略 $\pi$ 是 $L_\pi$-LC 的，若 $\forall s,s'\in\statespace$
    \begin{align*}
        \wassersteindistance[1](\pi(\cdot\vert s)\Vert \pi(\cdot\vert s'))\leq L_\pi \distance_{\statespace}(s,s').
    \end{align*}
\end{defi}

这些假设具有重要意义。例如，通过保证某个动作在状态空间的球域内占优，可以简化对主导动作的识别 \cite{rachelson2010locality}。同一文献还提供了关于 Q 函数 Lipschitz 性质的结果，如下所述。

\begin{thm}[Q函数的Lipschitz性质，\cite{rachelson2010locality} 的定理 1]
\label{th:lip_Q_function}
    考虑一个 $(L_P,L_r)$-LC 的 \gls{mdp} 和一个 $L_\pi$-LC 的策略 $\policy$。若 $\gamma L_P(1+L_\policy)\le 1$，则 $Q^\pi$ 是 $L_Q$-LC 的，其中
    \begin{align*}
        L_Q=\frac{L_r}{1-\gamma L_P(1+L_\policy)}.
    \end{align*}
\end{thm}

最后，\cite{metelli2020control} 引入了关于时间的Lipschitz性质的更近期概念（假设 4.1）。该假设在延迟情况下尤其有用，因为轨迹的光滑性是智能体预测其动作效果的关键因素。

\begin{defi}[时间Lipschitz MDP，\cite{metelli2020control} 的假设 4.1]
\label{def:time_lip}
    若 $\forall (s,a)\in \statespace\times\actionspace$ 有
    \begin{align*}
        \wassersteindistance[1](\transitionfunction(\cdot\vert s,a)\Vert \delta_s)\le L_T,
    \end{align*}
    其中 $\delta_s$ 由Dirac测度 $\delta$ 定义为 $\delta_s(x)=\delta(x-s)$，则称该 \gls{mdp} 是 $L_T$-时间Lipschitz连续的（$L_T$-TLC）。
\end{defi}

该定义看似与传统的Lipschitz性质定义不一致，但本文将展示它可视为关于已历经步数的Lipschitz性质。实际上，右端可理解为 $L_T\cdot 1$，其中 $1$ 对应于 \gls{mdp} 的时间单位（即步）。对于左端，可将 $p(\cdot\vert s,a)$ 重写为 $p^1(\cdot\vert s,a)$，以强调转移发生在一步之内。考虑执行多个动作的情况，例如 $n$ 个动作的序列 $\boldsymbol a=(a_1,\dots,a_n)$，可引入符号 $p^n(\cdot\vert s,\boldsymbol a)$ 表示从状态 $s$ 开始顺序执行 $\boldsymbol a$ 中的动作。进一步可将符号扩展为 $p^0(\cdot\vert s)=\delta_s$，表示在状态 $s$ 不执行任何动作的效果。采用这一新符号后，TLC 定义中的等式变为：
\begin{align*}
    \wassersteindistance[1](p^1(\cdot\vert s,a)\Vert p^0(\cdot\vert s))\le L_T\cdot (1-0).
\end{align*}
此时更清晰地表明Lipschitz性质是针对 \gls{mdp} 的步数而言的。将在 \Cref{pp:tlc_delay} 中证明 $p_{\boldsymbol a}^n(\cdot\vert s)$ 与 $p^0(\cdot\vert s)$ 之间距离的推广结果。


\subsection{POMDP}
\label{subsec:pomdp}

本小节简要介绍 \gls{pomdp}（部分可观测马尔可夫决策过程）\cite{kaelbling1998planning} 的概念，后文将看到它与延迟过程具有相似性。\gls{pomdp} 扩展了 \gls{mdp} 框架，用于处理智能体无法观测当前状态、只能获取部分信息的情形。更形式化地，\gls{pomdp} 是一个七元组 $(\statespace,\actionspace, \transitionfunction, \rewardfunction, \initialstatedistribution, \Omega, O)$，与 \gls{mdp} 相比增加了两个元素。首先，$\Omega$ 是智能体可能接收到的观测集（而非状态本身）。其次，$O:\statespace\times\actionspace\mapsto \setofprobability(\Omega)$ 称为观测函数，定义了 $\Omega$ 上的概率分布。具体而言，若智能体通过执行动作 $a$ 到达状态 $s'$，则 $O(o; s',a)$ 是智能体收到观测 $o$ 的概率。奖励函数和转移函数仍依赖于环境的当前状态，但该状态对智能体是未知的。优化目标仍为 \Cref{subsec:objective} 所定义的目标。

求解 \gls{pomdp} 的传统方法是维护当前状态的概率分布（称为置信状态），并通过每次收集到的新观测更新该概率。如此可将问题转化回 \gls{mdp}，其中状态即为上述置信状态。类似概念在控制理论文献 \cite{kumar2015stochastic} 中称为信息状态（information state），后文将在 \Cref{sec:related_works} 中看到，该概念也被用于延迟问题。关于 \gls{pomdp} 的更多信息，请参阅 \cite{hauskrecht2000value, spaan2012partially}。
